%--------------------------------------------------------------------------
% CONTRIBUCIONES ACADÉMICAS DEL AUTOR
%--------------------------------------------------------------------------
\chapter*{Síntesis de Aportes del Proyecto}
\addcontentsline{toc}{chapter}{Síntesis de Aportes del Proyecto}

Este trabajo de grado aporta una solución aplicada a un problema real de gestión documental y trazabilidad operativa en una empresa contratista multiservicio. Las contribuciones se organizan en frentes académicos, técnicos y metodológicos, evitando presentar como resultado operativo aquello que no cuente con evidencia verificable anexada.

\section*{Ingeniería de software y arquitectura}
\begin{itemize}
  \item Definición de una arquitectura web modular con separación de responsabilidades entre presentación, lógica de negocio y persistencia documental.
  \item Consolidación de contratos compartidos para reducir inconsistencias entre \frontend{} y \backend{}.
  \item Organización de la autenticación y la autorización bajo un esquema de control de acceso basado en roles (RBAC).
  \item Documentación de decisiones arquitectónicas relacionadas con la organización del monorepo, la persistencia documental, la autenticación y la operación con conectividad variable.
\end{itemize}

\section*{Gestión documental y trazabilidad}
\begin{itemize}
  \item Estructuración de un flujo documental orientado a órdenes de trabajo, evidencias, informes técnicos, actas y soportes de cierre.
  \item Relación explícita entre registros operativos, documentos técnicos y seguimiento administrativo de SES, facturación y pago.
  \item Planteamiento de una línea futura para formularios digitales reutilizables a partir de documentos heredados, sujeta a aprobación formal y validación posterior.
\end{itemize}

\section*{Operación en campo y continuidad del servicio}
\begin{itemize}
  \item Diseño de un enfoque de captura de información en campo para contextos de conectividad limitada.
  \item Integración conceptual de mecanismos de sincronización y persistencia local orientados a continuidad operativa.
  \item Enfoque en evidencias fotográficas y documentación de actividades para fortalecer la trazabilidad del servicio.
\end{itemize}

\section*{Contribuciones metodológicas}
\begin{itemize}
  \item Aplicación de una metodología orientada a diagnóstico, diseño, implementación y validación técnica de un artefacto de software en contexto empresarial.
  \item Definición de un flujo de trabajo \textit{contract-first}, desde esquemas de validación hasta componentes de interfaz, para mantener trazabilidad entre datos y experiencia de usuario.
  \item Propuesta de compuertas de calidad por iteración, incluyendo revisión de tipos, pruebas automatizadas, construcción del proyecto y verificación funcional.
  \item Organización de historias de usuario, criterios de aceptación y matriz de trazabilidad entre objetivos, módulos, pruebas y evidencias.
\end{itemize}

\section*{Contribuciones al estado del arte aplicado}
\begin{itemize}
  \item Revisión de plataformas comerciales de gestión de servicios de campo y mantenimiento, con atención a su aplicabilidad para pequeñas y medianas empresas contratistas.
  \item Análisis comparativo de repositorios \textit{open source} de CMMS, FSM y ERP, identificando patrones arquitectónicos y brechas funcionales frente al contexto de CERMONT S.A.S.
  \item Evaluación de librerías de formularios dinámicos y herramientas de extracción documental como base para una evolución futura del sistema.
  \item Elaboración de criterios para justificar el desarrollo de una solución a medida frente a la adopción directa de software comercial o genérico.
\end{itemize}

\section*{Formación académica y transferencia}
\begin{itemize}
  \item Articulación entre práctica empresarial, ingeniería electrónica, desarrollo de software y transformación digital aplicada.
  \item Producción de una referencia técnica reutilizable para proyectos similares en organizaciones contratistas multiservicio.
  \item Fortalecimiento de competencias en análisis de procesos, arquitectura de software, validación y documentación académica de soluciones tecnológicas.
\end{itemize}
